% Generated draft: requires mathematical and source review.
\documentclass[11pt]{article}
\usepackage[margin=25mm]{geometry}
\usepackage{fontspec}
\setmainfont{Latin Modern Roman}
\usepackage{amsmath,amssymb,mathtools,unicode-math}
\setmathfont{Latin Modern Math}
\usepackage{enumitem,needspace,xurl,hyperref}
\setlength{\parindent}{0pt}
\setlength{\parskip}{5pt}
\setlength{\emergencystretch}{3em}
\newcommand{\N}{\mathbb{N}}
\newcommand{\Z}{\mathbb{Z}}
\newcommand{\Q}{\mathbb{Q}}
\newcommand{\R}{\mathbb{R}}
\newcommand{\C}{\mathbb{C}}
\newcommand{\F}{\mathbb{F}}
\newcommand{\A}{\mathbb{A}}
\newcommand{\PP}{\mathbb{P}}
\newcommand{\E}{\mathbb{E}}
\newcommand{\eps}{\varepsilon}
\newcommand{\dd}{\,\mathrm{d}}
\begin{document}
\section*{UnsolvedMath ID: 701269}
\textbf{Source problem code:} catalogue entry
\subsection{Definitions and mathematical statement}
Let \(F(x,y)\) be the free group on two generators. A sequence of words is \(u=(u_n(x,y))_{n\geq 1}\) with \(u_n\in F(x,y)\); its evaluation at \((a,g)\in G^2\) is denoted by \(u_n(a,g)\). For such a sequence, define
\[
E_u(G)=\{g\in G:\exists n\geq 1\ \forall a\in G,\ u_n(a,g)=1_G\}.
\]
These are the \(u\)-Engel elements. The index \(n\) may depend on \(g\) in the given group \(G\), but must work simultaneously for every \(a\in G\).

For a group \(H\), put \(H^{(0)}=H\) and \(H^{(j+1)}=[H^{(j)},H^{(j)}]\), where \([K,K]\) is the subgroup generated by all \(a^{-1}b^{-1}ab\) with \(a,b\in K\). The group \(H\) is solvable if \(H^{(m)}=\{1\}\) for some integer \(m\geq 0\). For a finite group \(G\), its solvable radical \(R_{\mathrm{sol}}(G)\) is its largest solvable normal subgroup.

The source requires the word sequence to satisfy the conditions of Kourovka Notebook Problem 15.75. Those conditions are not included in the supplied record. Write \(\mathcal U_{15.75}\) provisionally for the class of sequences satisfying that unavailable requirement; this notation does not supply its missing definition.

The question attributed to B. I. Plotkin is whether
\[
\exists u\in\mathcal U_{15.75}\ \forall\text{ finite groups }G,
\qquad R_{\mathrm{sol}}(G)=E_u(G).
\]
Equivalently, one asks for a single admissible sequence, independent of \(G\), such that for every finite group \(G\) and every \(g\in G\),
\[
g\in R_{\mathrm{sol}}(G)
\quad\Longleftrightarrow\quad
\exists n\geq 1\ \forall a\in G,\ u_n(a,g)=1_G.
\]
The admissibility requirement from Problem 15.75 is essential to the stated question; without it, this is only a conditional formulation, not a complete standalone statement.

As motivation, the supplied source states that in every Noetherian group the nilpotent radical coincides with the collection of all Engel elements, citing R. Baer and B. I. Plotkin. Here Noetherian means that every ascending chain of subgroups stabilizes. This background assertion is reported from the source and is not independently verified here.
\subsection{Short English statement}
Does there exist a single sequence of two-variable group words, satisfying the unspecified conditions of Kourovka Notebook Problem 15.75, that detects membership in the solvable radical of every finite group? For each element, one word in the sequence must vanish for every choice of the other variable. The missing cross-reference prevents a complete standalone formulation.
\subsection{Sources}
\begin{itemize}
\item Local UnsolvedMath registry, supplied complete source record: permanent ID \(701269\), exact source code \texttt{catalogue entry}, title \emph{Kourovka Notebook Problem 16.77}. The record labels the problem open; this is a supplied status, not an independently verified literature assessment.
\item \emph{The Kourovka Notebook (Unsolved Problems in Group Theory)}, Problem 16.77, proposed by B. I. Plotkin in the 16th issue (2006), as represented in the supplied reference to the 21st issue edition (2026), \texttt{arXiv:1401.0300v48}. The supplied collection description names E. I. Khukhro and V. D. Mazurov as editors.
\item Kourovka Notebook Problem 15.75: referenced by the statement for the required conditions on the word sequence, but its text is absent from the supplied record.
\item R. Baer, \emph{Math. Ann.}, \textbf{133} (1957), 256--270. Supplied citation supporting the motivational assertion about Noetherian groups; no article title is supplied.
\item B. I. Plotkin, \emph{Izv. Vysš. Učebn. Zaved. Mat.}, \textbf{1958}, no. 1(2), 130--135 (Russian). Supplied citation supporting the same motivational assertion; no article title is supplied.
\end{itemize}
\end{document}
